% META: {"id": "1NSI-LANG-M3", "chapitre": "1NSI-LANGAGE", "type_objet": "methode", "capacites": ["P-LANG-05"], "capacites_codes": ["C5"], "methodes": ["M3"], "mode_creation": "ex_nihilo", "sources_inspiration": [], "status": "generated"}
\begin{fichemethode}{M3}{Trouver une fonction dans la documentation d'une bibliothèque}
\textbf{Quand l'utiliser ?} Quand tu as besoin d'une opération que tu ne veux pas
reprogrammer — trier, tirer au hasard, lire une date — et que tu supposes qu'elle existe
déjà. Aucune connaissance exhaustive des bibliothèques n'est exigée : savoir chercher
remplace savoir par cœur.

\textbf{Pas à pas :}
\begin{enumerate}
  \item \textbf{Nomme ton besoin en une phrase} : « tirer un entier au hasard entre deux
    bornes ». Les mots de cette phrase sont tes mots-clés.
  \item \textbf{Devine le module.} Le hasard, c'est \lstinline{random} ; les maths,
    \lstinline{math} ; les dates, \lstinline{datetime} ; les fichiers tabulés,
    \lstinline{csv}. En cas de doute, \lstinline{help("modules")} les liste.
  \item \textbf{Liste ce que le module offre} avec \lstinline{dir(module)}, et repère les
    noms proches de ton besoin.
  \item \textbf{Lis la fiche de la fonction} avec \lstinline{help(module.fonction)}. Trois
    choses seulement t'intéressent : les \textbf{paramètres} et leur ordre, la
    \textbf{valeur de retour}, et les \textbf{bornes} — incluses ou exclues ?
  \item \textbf{Essaie sur un cas dont tu connais la réponse} avant de l'employer pour de
    bon. Une fonction mal comprise est pire qu'une fonction absente : elle ne plante pas.
\end{enumerate}

\textbf{Exemple :} « tirer un entier au hasard entre $1$ et $6$ ».

\begin{python}
import random

[n for n in dir(random) if "int" in n]
# ['getrandbits', 'randint', 'randrange']
help(random.randint)
# randint(a, b) -> Return random integer in range [a, b],
#                  including both end points.
\end{python}

La fiche dit \textbf{including both end points} : les deux bornes sont possibles.
\lstinline{random.randint(1, 6)} convient. Attention, \lstinline{random.randrange(1, 6)}
n'atteint jamais $6$ : la borne haute y est exclue, comme dans \lstinline{range}.

% BEGIN-VERIFY
% import random
% # randint : les deux bornes sont atteignables
% random.seed(2026)
% tirages = [random.randint(1, 6) for _ in range(2000)]
% assert min(tirages) == 1 and max(tirages) == 6
% assert set(tirages) == {1, 2, 3, 4, 5, 6}
% # randrange : la borne haute est exclue, comme range
% random.seed(2026)
% autres = [random.randrange(1, 6) for _ in range(2000)]
% assert max(autres) == 5 and 6 not in autres
% # la docstring dit exactement cela, et c'est elle qui tranche
% assert "including both end points" in random.randint.__doc__
% # le module expose bien ces noms : dir() permet de les trouver
% assert "randint" in dir(random) and "randrange" in dir(random)
% # une graine fixee rend le tirage reproductible : indispensable pour tester
% random.seed(1); premier = random.randint(1, 6)
% random.seed(1); assert random.randint(1, 6) == premier
% END-VERIFY

\textbf{Pièges :}
\begin{itemize}
  \item \textbf{Supposer les bornes.} \lstinline{randint(1, 6)} atteint $6$,
    \lstinline{randrange(1, 6)} ne l'atteint jamais, \lstinline{range(1, 6)} non plus. Un
    dé programmé avec \lstinline{randrange} ne fait jamais $6$, et le programme ne signale
    rien.
  \item Recopier un exemple trouvé en ligne sans lire la fiche : l'ordre des paramètres ou
    la valeur de retour peuvent différer d'une version à l'autre.
  \item Réécrire une fonction qui existe déjà : plus longue à écrire, et plus sûrement
    fausse.
  \item Ne pas fixer de graine quand on teste du hasard : le test passe une fois sur deux,
    et personne ne sait pourquoi. \lstinline{random.seed(...)} rend le tirage reproductible.
\end{itemize}

\textbf{Vérifier :} tire mille valeurs et regarde le minimum et le maximum obtenus. Ils
doivent être exactement les bornes que tu attends. C'est le contrôle qui démasque une borne
exclue prise pour incluse.

\textbf{S'entraîner :} exercices C5 du chapitre.
\end{fichemethode}
